% from ega3-7-fr.tex
\oldpage[III]{184}

\begin{env}[7.3.8]
Si $A$ est un anneau commutatif, $B$ une $A$-algèbre (non nécessairement
commutative), $T_\bullet$ un foncteur homologique covariant de
$\mathbf{Ab}_A$ dans $\mathbf{Ab}$, il résulte des définitions (7.1.3) que le
foncteur de $\mathbf{Ab}_B$ dans $\mathbf{Ab}$ obtenu par extension des
scalaires de $A$ à $B$, et que nous noterons
\[
  T_\bullet^{(B)}=(T_i^{(B)}),
\]
est encore un foncteur homologique.
\end{env}

\begin{corollary}[7.3.9]
Supposons que $T_\bullet$ vérifie les conditions générales de (7.3.7) et
commute aux limites inductives, et en outre que $A$ soit un anneau intègre et
tous les $T_n(A)$ des $A$-modules de présentation finie. Alors, pour tout
entier $N$, il existe un $f\in A-\{0\}$ tel que le foncteur
\[
  T_p^{(A_f)}:\mathbf{Ab}_{A_f}\longrightarrow\mathbf{Ab}
\]
soit exact pour $p\leq N$.
\end{corollary}

\begin{proof}
Par hypothèse, $T_i$ est exact pour $i\leq i_0$, donc $T_i(A)$ est plat pour
ces valeurs de $i$. En vertu de (7.3.7, $b)$), il suffit de prendre $f$ tel que
\[
  T_p^{(A_f)}(A_f)=T_p(A_f)
\]
soit un $A_f$-module libre pour $i_0<p\leq N$. Or, on a
$T_p(A_f)=(T_p(A))_f$ puisque $T_p$ commute aux limites inductives (7.1.4). Si
$x$ est le point générique de $\operatorname{Spec}(A)$, $(T_p(A))_x$ est un
espace vectoriel de dimension finie sur le corps des fractions de $A$. Comme
chaque $T_p(A)$ est de présentation finie, il existe bien un $f$ ayant la
propriété voulue (Bourbaki, \emph{Alg. comm.}, chap. II, \S 5,
n\textsuperscript{o} 1, cor. de la prop. 2).

On notera que s'il n'y a qu'un nombre fini d'indices $i$ tels que $T_i\ne0$,
il existe $f\in A-\{0\}$ tel que tous les $T_p^{(A_f)}$ soient exacts.
\end{proof}

\begin{corollary}[7.3.10]
Supposons que $T_\bullet$ vérifie les conditions générales de (7.3.7) et
commute aux limites inductives, que $A$ soit commutatif et noethérien et les
$T_n(A)$ des $A$-modules de type fini. Alors, pour tout entier $N$, il existe
un ouvert dense $U$ de $\operatorname{Spec}(A)$ tel que, pour tout $p\leq N$,
la fonction
\[
  x\mapsto\operatorname{rang}_{\kappa(x)}(T_p(\kappa(x)))
\]
soit constante dans $U$.
\end{corollary}

\begin{proof}
Soit $\mathfrak p$ un idéal premier minimal de $A$; par hypothèse, l'anneau
$B=A/\mathfrak p$ est intègre et $\operatorname{Spec}(B)$ s'identifie à une
composante irréductible de l'espace topologique $\operatorname{Spec}(A)$. Nous
allons montrer, par récurrence sur $p\leq N$ qu'il existe
$f_p\in B-\{0\}$ tel que si on pose $B'=B_{f_p}$, $T_i^{(B')}$ soit exact et
les $T_i(B')$ des $B'$-modules de type fini pour $i\leq p$. La proposition est
vraie pour $p\leq i_0$ en vertu de l'hypothèse, en prenant $f_p=1$ (donc
$B'=B=A/\mathfrak p$) car $T_p$ étant alors exact, $T_p(B)$ est isomorphe à
$T_p(A)/T_p(\mathfrak p)$, donc est un $A$-module (et \emph{a fortiori} un
$B$-module) de type fini. Raisonnons par récurrence sur $p$; $f_p$ est l'image
canonique dans $B$ d'un élément $g_p\in A$, et si l'on pose $A'=A_{g_p}$, on a
$B'=A'/\mathfrak p'$ où $\mathfrak p'$ est un idéal premier minimal de $A'$,
égal à $\mathfrak p_{g_p}$. Comme on a
\[
  T_i(A_{g_p})=(T_i(A))_{g_p},
\]
les $T_i(A')$ sont des $A'$-modules de type fini, donc le foncteur
$T_\bullet^{(A')}$ vérifie les mêmes hypothèses que $T_\bullet$, mais en
remplaçant $i_0$ par $p$. On peut donc se borner au cas où $A'=A$, et où $T_p$
est exact; la suite exacte
\[
  0\longrightarrow\mathfrak p\longrightarrow A\longrightarrow A/\mathfrak p
  \longrightarrow0
\]
donne alors la suite exacte
\[
  T_{p+1}(A)\longrightarrow T_{p+1}(A/\mathfrak p)
  \xrightarrow{\partial}T_p(\mathfrak p)\longrightarrow T_p(A),
\]
et comme $T_p$ est exact, la dernière flèche de droite est injective, donc
$T_{p+1}(A/\mathfrak p)$ est un quotient de $T_{p+1}(A)$ et est par suite de
type fini. Notons maintenant que le raisonnement de (7.3.9) n'a utilisé le fait
que les $T_p(A)$ sont de type fini que pour $p\leq N$; on peut donc l'appliquer
à l'anneau intègre $B$, au foncteur $T_\bullet^{(B)}$ et à $N=p+1$, ce qui
achève le raisonnement par récurrence. Cela étant, il y a un
$f_N\in B-\{0\}$ tel que $\operatorname{Spec}(B_{f_N})$ soit un ouvert $V$
partout dense dans $\operatorname{Spec}(B)$ ne rencontrant aucune autre
composante irréductible de $\operatorname{Spec}(A)$. Si la proposition est
